\documentclass{report}
 \usepackage{newcent}
 \usepackage{pifont}
 \usepackage[dvips]{color}

 \usepackage{graphics}

 \parindent 0pt
 \pagestyle{plain}
 \topmargin0cm
 \headheight0pt
 \headsep0pt
 \textwidth16.5cm
 \textheight23cm
 \footskip1cm
 \oddsidemargin0cm
 \evensidemargin0cm

 \renewcommand{\labelitemi}{\ding{228}}
 \renewcommand{\labelitemii}{\ding{169}}
 \renewcommand{\labelitemiii}{\ding{72}}
 \renewcommand{\arraystretch}{1.1}
 \begin{document}
 \vskip \baselineskip
 \vbox{\hsize \textwidth
 \hrule
 \vbox{\rule[-.55em]{0pt}{1.7em}\strut \vrule \quad \hfill
 \label{S1Section}{\bf S1 Section } \hfill
 \quad \vrule}\hrule}
 \vskip \baselineskip
 {The first Section in the document. \hskip1pt\\
 Some more paragraphs... }
 \vskip \baselineskip
 \vbox{\hsize \textwidth
 \hrule
 \vbox{\rule[-.55em]{0pt}{1.7em}\strut \vrule \quad 
 \label{S2Subsection}{\bf S2 Subsection } \hfill 
 {\it  S1 Section}\quad \vrule}\hrule}
 \vskip \baselineskip
 {The first subsection inside the section. }
 \vskip \baselineskip
 \vbox{\hsize \textwidth
 \hrule
 \vbox{\rule[-.55em]{0pt}{1.7em}\strut \vrule \quad 
 \label{S2Subsection2}{\bf S2 Subsection 2 } \hfill 
 {\it  S1 Section}\quad \vrule}\hrule}
 \vskip \baselineskip
 {The next subsection. }
 \vskip \baselineskip
 \vbox{\hsize \textwidth
 \hrule
 \vbox{\rule[-.55em]{0pt}{1.7em}\strut \vrule \quad 
 \label{S3Unit}{\bf S3 Unit } \hfill
 {\it  S2 Subsection 2}\quad \vrule}\hrule}
 \vskip \baselineskip
 {This kind of unit is often used. }\vskip .5\baselineskip 
 {\parindent 0pt \vbox{\bf \label{S4NamedParagraph}S4 Named Paragraph }
 \leftskip 20pt \rightskip 20pt \vskip .5\baselineskip\begin{list}{}{}\item{
 \parindent 0pt 
 {More words... }}\end{list}}
 \vskip .5\baselineskip
 \vskip 2\baselineskip
 \vbox{\hfil
 \vbox{\hrule
 \hbox{\vrule\qquad
 {\rule[-.55em]{0pt}{1.7em}\sc Example}\quad\vrule}
 \hrule}
 }
 \nobreak\vskip \baselineskip\nobreak
 \halign to \textwidth{
 \textcolor{black}{\strut\vrule width 5pt#}\hskip 2em& 
   #\hfil  \cr
  &\cr 
 & {\tt \hskip 1em \hfil} \cr
 & {\tt \hskip 1em open Amoeba \hfil} \cr
 & {\tt \hskip 1em open Stderr \hfil} \cr
 & {\tt \hskip 1em open Unix \hfil} \cr
 & {\tt \hskip 1em let $\_$ $=$ \hfil} \cr
 & {\tt \hskip 1em {\tt\hskip 2em}let a $=$ 1 in \hfil} \cr
 & {\tt \hskip 1em {\tt\hskip 2em}let b $=$ \hfil} \cr
 & {\tt \hskip 1em {\tt\hskip 2em}{\tt\hskip 2em}let c $=$ 1 in \hfil} \cr
 & {\tt \hskip 1em {\tt\hskip 2em}\hfil} \cr }
 \vskip 2\baselineskip
 {A data format table: }
 \vskip 2\baselineskip
 \begin{center}\vbox{
 {\bf FLIP protocol header format }\\\begin{tabular}
 {|p{3em}|p{3em}|p{3em}|p{3em}|p{3em}|p{3em}|p{3em}|p{3em}|}\hline
 {\bf 7 }&
 {\bf 6 }&
 {\bf 5 }&
 {\bf 4 }&
 {\bf 3 }&
 {\bf 2 }&
 {\bf 1 }&
 {\bf 0 }
 \\\hline\hline
 \multicolumn{2}{|l|}{Max HopCnt }&\multicolumn{2}{l|}{Actual HopCnt }&Res. &
 Res. &Type &Vers. \\\hline
 \end{tabular}}
 \end{center}
 \vskip 2\baselineskip
 {The following programming interface gives an impression of the MANDoc system. 
 }
 \vskip 2\baselineskip
 \vbox{\hfil
 \vbox{\hrule
 \hbox{\vrule\qquad
 {\rule[-.55em]{0pt}{1.7em}\sc Programming Interface: {\sl ML }}\quad\vrule}
 \hrule}
 }
 \nobreak\vskip \baselineskip\nobreak
 {
 \nobreak\vskip .5\baselineskip 
 \halign{
 \textcolor{black}{\strut\vrule width 5pt\hskip 20pt \hfill #} & # \hfil & # & #\hfill & # \hfil & \qquad # \hfill \cr
 &&&&& \cr
 $[$ & {\it ret1 }: int $*$ & & & & \cr
 & {\it ret2 }: bool  & $]=$ & {\bf ml$\_$fun } & 
  $^\sim$arg : float $\rightarrow$ & \cr
 & & & &  $^\sim$kind : bool  & \cr
 }
 \halign{
 \textcolor{black}{\strut\vrule width 5pt\hskip 20pt \hfill } # & # \hfil  & \qquad # \cr
 &&\cr
 \textcolor{black}{external} {\bf buf$\_$physical }: 
 & size :  int  $\rightarrow$  &  \cr
 &  buffer  &  \cr
  & 
 {\bf $=$ "ext$\_$buf$\_$physical" } & \cr
 }\halign{\textcolor{black}{\strut\vrule width 5pt#}&#\cr&\cr}
 }
 \vskip 2\baselineskip
 {Not only ML, C interfaces are supported, too. }
 \vskip 2\baselineskip
 \vbox{\hfil
 \vbox{\hrule
 \hbox{\vrule\qquad
 {\rule[-.55em]{0pt}{1.7em}\sc Programming Interface: {\sl ML }}\quad\vrule}
 \hrule}
 }
 \nobreak\vskip \baselineskip\nobreak
 {
 \nobreak\vskip .5\baselineskip 
 \halign{
 \textcolor{black}{\strut\vrule width 5pt\hskip 20pt \hfill #} & # \hfil & # & #\hfill & \qquad # \hfill \cr
 &&&&\cr
 &struct module$\_$s  & {\bf mycfun }$($&int a ,  & \cr
 & & &float b ,  & \cr
 & & &struct module$\_$s $*$p $);$ & \cr
 }
 \halign{
 \textcolor{black}{\strut\vrule width 5pt\hskip 20pt \hfill #} & # & #\cr
 &&\cr
 &struct model  & {\bf myvar };\cr
 }
 \halign{
 \textcolor{black}{\strut\vrule width 5pt\hskip 20pt \hfill #} & # \hfil & # \hfil & #\hfill & \qquad # \hfill \cr
 &&&&\cr
 \textcolor{black}{struct} & {\bf c$\_$struc } $\{$ & &char name[20] $;$ & \cr
 && &int len $;$ & \cr
 && &float zahl  $\}$  & \cr
 }\halign{\textcolor{black}{\strut\vrule width 5pt#}&#\cr&\cr}
 }
 \vskip 2\baselineskip
 {Now we have an image: }
 \vskip \baselineskip
 \begin{center}\begin{minipage}{.9\textwidth}\vbox{
 \hbox{\includegraphics{logo.eps}}
 {(Fig. 1) All the components of the VAMNET system together in an example 
 configuration. \hskip1pt\\
 }}\end{minipage}\end{center}
 \vskip \baselineskip
 \vskip \baselineskip
 \vbox{\hsize \textwidth
 \hrule
 \vbox{\rule[-.55em]{0pt}{1.7em}\strut \vrule \quad 
 \label{S2single}{\bf S2 single } \hfill 
 \quad \vrule}\hrule}
 \vskip \baselineskip
 {One more isolated S2 section. A table: }
 \vskip 2\baselineskip
 \begin{center}
 \begin{tabular}{|p{6.6cm}|p{1.8cm}|}\hline
 \multicolumn{2}{|c|}{{\sc Table }}\\\hline\hline
 The special interfaces are used to show ML or C functions, types, submodules 
 and values. Only the content must be given by the user, not the alignment. 
 Commonly, they are collected in an interface subunit. &1+1$=$2 \\\hline
 \end{tabular}
 \end{center}
 \vskip 2\baselineskip
 \vskip 2\baselineskip
 \begin{center}
 \begin{tabular}{|p{3cm}|p{1.8cm}|}\hline
 \multicolumn{2}{|c|}{{\sc Table }}\\\hline\hline
 Operation &1+1$=$2 \\\hline
 \end{tabular}
 \end{center}
 \vskip 2\baselineskip
 \vskip \baselineskip
 \vbox{\hsize \textwidth
 \hrule
 \vbox{\rule[-.55em]{0pt}{1.7em}\strut \vrule \quad \hfill
 \label{NextSection1}{\bf Next Section 1 } \hfill
 \quad \vrule}\hrule}
 \vskip \baselineskip
 {Here comes another section following the above main section. }
 \end{document}
 